% ============================================================
%  SCHOLA DOMESTICA — Magister Mathematica
%  Level 1 | Week 10 | Day 4
%  Topic: All Addition Facts — Sums to 8, Mixed Review
% ============================================================
\documentclass[12pt, letterpaper]{article}

\usepackage[top=1in, bottom=1in, left=1in, right=1in]{geometry}
\usepackage[T1]{fontenc}
\usepackage[utf8]{inputenc}
\usepackage{lmodern}
\usepackage{microtype}
\usepackage{amsmath}
\usepackage{xcolor}
\usepackage{tikz}
\usepackage{array}
\usepackage{enumitem}
\usepackage{multicol}
\usepackage{mdframed}
\usepackage{fancyhdr}

\definecolor{scholared}{RGB}{160,30,30}
\definecolor{scholagold}{RGB}{180,140,40}
\definecolor{scholacream}{RGB}{253,248,235}
\definecolor{scholagreen}{RGB}{50,110,60}
\definecolor{scholablue}{RGB}{40,80,140}

\pagestyle{fancy}
\fancyhf{}
\renewcommand{\headrulewidth}{0.6pt}
\fancyhead[L]{\small\textcolor{scholared}{\textbf{Schola Domestica — Magister Mathematica}}}
\fancyhead[R]{\small\textcolor{scholared}{Level 1 · Week 10 · Day 4}}
\fancyfoot[C]{\small\textcolor{gray}{— \thepage\ —}}

\newcommand{\answerline}{\underline{\hspace{2.5cm}}}
\newcommand{\biganswerline}{\underline{\hspace{4cm}}}
\newcommand{\numbox}{\tikz{\draw[scholablue,thick](0,0)rectangle(1.2cm,0.85cm);}}

\newmdenv[
  backgroundcolor=scholacream,
  linecolor=scholared,
  linewidth=1.5pt,
  roundcorner=6pt,
  innertopmargin=8pt,
  innerbottommargin=8pt,
  innerleftmargin=10pt,
  innerrightmargin=10pt
]{lessonbox}

\newmdenv[
  backgroundcolor={rgb,255:red,235;green,245;blue,235},
  linecolor=scholagreen,
  linewidth=1.2pt,
  roundcorner=4pt,
  innertopmargin=6pt,
  innerbottommargin=6pt,
  innerleftmargin=8pt,
  innerrightmargin=8pt
]{enrichbox}

\begin{document}

\begin{center}
  {\LARGE\textbf{\textcolor{scholared}{All Facts to Eight:}}}\\[4pt]
  {\Large\textcolor{scholagold}{\textit{Putting It All Together}}}\\[6pt]
  {\small Name: \biganswerline \hspace{1cm} Date: \answerline}
\end{center}

\vspace{4pt}
\noindent\textcolor{scholared}{\rule{\linewidth}{1pt}}
\vspace{4pt}

% IMAGE PROMPT (Beatrix Potter watercolour style):
% A badger sits at a little wooden desk with an abacus, carefully pushing beads. Autumn leaves
% visible through a round window. Earthy warm tones, white paper background, no text.

\begin{lessonbox}
\textbf{Your growing toolkit of strategies:}
\begin{itemize}[leftmargin=1.5em, itemsep=2pt]
  \item \textbf{$+0$}: same number \quad \textbf{$+1$}: one more \quad \textbf{$+2$/$+3$}: count on
  \item \textbf{$+4$/$+5$}: count on from the bigger addend
  \item \textbf{Doubles}: memorise 1+1 through 4+4
  \item \textbf{Near-doubles}: use the double, then $\pm 1$
\end{itemize}
\end{lessonbox}

\bigskip

% ===== PART I: HORIZONTAL FACTS =====
\noindent\textcolor{scholared}{\large\textbf{Part I — Quick Facts}} \hfill \textit{(Problems 1–8)}

\medskip

\noindent
\begin{tabular}{@{}p{3.3cm} p{3.3cm} p{3.3cm} p{3.3cm}@{}}
\textbf{1.}\; $6 + 2 =$ \numbox &
\textbf{2.}\; $4 + 4 =$ \numbox &
\textbf{3.}\; $5 + 3 =$ \numbox &
\textbf{4.}\; $7 + 1 =$ \numbox \\[12pt]
\textbf{5.}\; $3 + 4 =$ \numbox &
\textbf{6.}\; $8 + 0 =$ \numbox &
\textbf{7.}\; $2 + 5 =$ \numbox &
\textbf{8.}\; $1 + 6 =$ \numbox
\end{tabular}

\bigskip\bigskip

% ===== PART II: VERTICAL FACTS =====
\noindent\textcolor{scholared}{\large\textbf{Part II — Vertical Facts}} \hfill \textit{(Problems 9–12)}

\bigskip

\noindent
\begin{tabular}{@{}p{2.0cm} p{2.0cm} p{2.0cm} p{2.0cm}@{}}
\textbf{9.}
$\begin{array}{r} 5 \\ +2 \\ \hline \phantom{0} \end{array}$
&
\textbf{10.}
$\begin{array}{r} 4 \\ +3 \\ \hline \phantom{0} \end{array}$
&
\textbf{11.}
$\begin{array}{r} 6 \\ +1 \\ \hline \phantom{0} \end{array}$
&
\textbf{12.}
$\begin{array}{r} 4 \\ +4 \\ \hline \phantom{0} \end{array}$
\end{tabular}

\bigskip\bigskip

% ===== PART III: MISSING ADDEND =====
\noindent\textcolor{scholared}{\large\textbf{Part III — Find the Missing Part}} \hfill \textit{(Problems 13–16)}

\bigskip

\noindent
\begin{tabular}{@{}p{3.5cm} p{3.5cm} p{3.5cm} p{3.5cm}@{}}
\textbf{13.}\; $4 +$ \numbox $= 8$ &
\textbf{14.}\; $\numbox + 5 = 7$ &
\textbf{15.}\; $3 +$ \numbox $= 7$ &
\textbf{16.}\; $\numbox + 4 = 6$
\end{tabular}

\bigskip\bigskip

% ===== PART IV: WORD PROBLEMS =====
\noindent\textcolor{scholared}{\large\textbf{Part IV — Story Problems}} \hfill \textit{(Problems 17–20)}

\medskip

\noindent\textbf{17.} A pond has \textbf{5} lily pads on the left and \textbf{3} on the right. How many lily pads in all?

\medskip\noindent\hspace{2em} $\numbox + \numbox = \numbox$ lily pads

\bigskip

\noindent\textbf{18.} A basket has \textbf{4} apples. Someone puts in \textbf{4} more. How many apples now?

\medskip\noindent\hspace{2em} $\numbox + \numbox = \numbox$ apples

\bigskip

\noindent\textbf{19.} A worm finds \textbf{2} leaves and some more leaves. It has \textbf{8} leaves in all. How many more did it find?

\medskip\noindent\hspace{2em} $2 +$ \numbox $= 8$\quad It found \numbox\ more leaves.

\bigskip

\noindent\textbf{20.} Write your own addition story that uses the number 7 as the whole:

\medskip\noindent\hspace{2em}\underline{\hspace{10cm}}

\medskip\noindent\hspace{2em} $\numbox + \numbox = 7$

\bigskip

\begin{enrichbox}
\textbf{\textcolor{scholagreen}{Oral Fact Drill — Sums to 8}}\\[3pt]
Teacher calls a fact; child answers immediately (no writing).\\[4pt]
$4{+}4$ \; $3{+}5$ \; $6{+}2$ \; $7{+}1$ \; $5{+}3$ \; $4{+}3$ \; $3{+}4$ \; $2{+}5$ \; $5{+}2$ \; $8{+}0$\\[2pt]
\textit{Goal: no hesitation. Repeat any that take more than 3 seconds.}
\end{enrichbox}

\vspace{0.3cm}
\noindent\textcolor{gray}{\small\textit{Ray's Primary Arithmetic — mastering all facts with sums to 8.}}

\end{document}
